\documentclass{article}

\usepackage{amsmath}
\usepackage{amsthm}
\usepackage{amssymb}
\usepackage[margin=1in]{geometry}

\newtheorem{theorem}{Theorem}
\newtheorem{lemma}{Lemma}
\newtheorem{corollary}{Corollary}
\newtheorem{proposition}{Proposition}
\theoremstyle{remark}
\newtheorem{remark}{Remark}

\newcommand{\Rcheck}{\check{R}}
\newcommand{\Ttilde}{\widetilde{T}}
\newcommand{\ord}{\operatorname{ord}}
\newcommand{\tr}{\operatorname{tr}}
\newcommand{\Tr}{\operatorname{Tr}}

\title{The Baxterised $\Omega$ is an Isaev--Kirillov Hecke $R$-factor staircase:\\ an exact spectral bridge}
\author{Clio}
\date{2026-05-30}

\begin{document}

\maketitle

\begin{abstract}
Clio's order law $\ord_{x=q^2} Z_\lambda(x) = \tau(\tau+1)/2$ governs a staircase monodromy $M(x)$ of Baxterised Hecke $R$-factors evaluated at a degenerate spectral point $x^* = q^2$. We prove that each such $R$-factor is, under an explicit M\"obius reparametrization of the spectral variable, \emph{exactly} the normalized Baxterised Hecke generator of Isaev--Kirillov (arXiv:1302.6495), whose Markov-traced staircases generate the Bethe commutative subalgebra of $H_q(S_n)$. Consequently the trace $Z_\lambda(x) = \tr_{V^\lambda} M(x)$ belongs to the Bethe-subalgebra trace family, and the order law is the vanishing order of a Bethe transfer-matrix trace at a fusion point. We record the bridge, its verification, and an honest account of what remains to identify $\Omega$ as a specific Bethe-subalgebra coefficient.
\end{abstract}

\section{Setup: Clio's object}

In the Hecke algebra $H_q(S_n)$ with generators $T_i$ satisfying
\[
(T_i - q)(T_i + q^{-1}) = 0,
\]
let
\[
P_q^{(i)} = \frac{T_i + 1}{q+1}, \qquad P_{-1}^{(i)} = \frac{qI - T_i}{q+1}
\]
be the projectors onto the $T_i = q$ and $T_i = -q^{-1}$ eigenspaces. (One checks $P_q^{(i)} + P_{-1}^{(i)} = \frac{(T_i+1) + (qI - T_i)}{q+1} = \frac{(q+1)I}{q+1} = I$; these are the idempotents I use throughout.) The Baxterised factor is
\[
\Rcheck_i(x) = P_q^{(i)} + r(x)\, P_{-1}^{(i)}, \qquad r(x) = \frac{q^2 - x}{q(x-1)}.
\]
The monodromy $M(x) = \prod_k \Rcheck_{i_k}(x)$ is the ordered product over a staircase reduced word; $Z_\lambda(x) = \tr_{V^\lambda} M(x, \dots, x)$ and $\Omega = (q+1)^N M(q^2)$.

The special values of $r$ organize the geometry of the spectral line:
\begin{itemize}
  \item $r(q) = 1$, so $\Rcheck(q) = I$: the \emph{regular} point;
  \item $r(q^2) = 0$, so $\Rcheck(q^2) = P_q$: the \emph{fusion} point onto the symmetric line;
  \item $r$ has a pole at $x = 1$: the other fusion, onto $P_{-1}$.
\end{itemize}
The order law, proved for hooks and intact-junction off-hooks, reads
\[
\ord_{x=q^2} Z_\lambda = \frac{\tau(\tau+1)}{2}, \qquad \tau = \max(0,\, 2\lambda'_1 - n - 1).
\]

The whole edifice is built from a single one-parameter family of operators $\Rcheck_i(x)$ on the Hecke algebra. It has long bothered me that this family looks like it ought to be \emph{somebody's} Baxterised generator. This note shows it is.

\section{The Isaev--Kirillov Baxterised generator and Bethe subalgebra}

Isaev--Kirillov (arXiv:1302.6495, Lett.\ Math.\ Phys.\ \textbf{104} (2014) 179) work on the same Hecke algebra. They write the quadratic relation as $T_i^2 = 1 + \lambda T_i$ with $\lambda := q - q^{-1}$, which is the same relation $(T_i - q)(T_i + q^{-1}) = 0$. They define the Baxterised generator
\[
T_i(x) = (1-x) T_i + \lambda x,
\]
with normalized form
\[
\Ttilde_i(x) = \frac{T_i(x)}{q - q^{-1} x}
\]
satisfying unitarity $\Ttilde_i(x)\, \Ttilde_i(x^{-1}) = I$ and the multiplicative-spectral Yang--Baxter equation.

They then build the Baxterised Jucys--Murphy element as a sandwiched staircase product
\[
y_n(x; \vec\xi) = T_{n-1}(x/\xi_{n-1}) \cdots T_1(x/\xi_1)\; y_1(x)\; T_1(x\xi_1) \cdots T_{n-1}(x\xi_{n-1}),
\]
and the Bethe commutative subalgebra $\mathcal{B}_n(\vec\xi) \subset H_q(S_n)$ is generated by the coefficients $\Phi_k$ of the Markov trace
\[
\tau_n(x) = \Tr_{(n+1)}\!\big( y_{n+1}(x; \vec\xi) \big) = \sum_k \Phi_k\, x^k.
\]
Commutativity $[\tau_n(x), \tau_n(z)] = 0$ is their Proposition 3.2. In the $q \to 1$ limit the $\Phi_k$ degenerate to the Gaudin Hamiltonians.

One bookkeeping difference must be flagged before any comparison. The IK spectral parameter is \emph{multiplicative}: their Yang--Baxter equation involves the product $xz$, not the sum $x + z$. My $r(x)$ is written \emph{additively}. This convention gap is exactly what the M\"obius bridge of the next section closes.

\section{Main Lemma: the spectral conjugacy}

All identities in this section were verified symbolically over $\mathbb{Q}(q)$ in sympy.

\begin{lemma}[Eigenvalues]
On the $T_i = q$ eigenline, $T_i(x)$ acts by $q - q^{-1} x$; on the $T_i = -q^{-1}$ eigenline, by $q x - q^{-1}$. Hence the normalized generator decomposes as $\Ttilde_i(x) = P_q + \rho(x) P_{-1}$ with eigenvalue-ratio
\[
\rho(x) = \frac{q x - q^{-1}}{q - q^{-1} x} = \frac{q^2 x - 1}{q^2 - x}.
\]
\end{lemma}

\begin{proof}
On the $T_i = q$ eigenline, $T_i(x) = (1-x) q + \lambda x = q - qx + (q - q^{-1})x = q - q^{-1}x$. On the $T_i = -q^{-1}$ eigenline, $T_i(x) = (1-x)(-q^{-1}) + \lambda x = -q^{-1} + q^{-1}x + (q - q^{-1})x = -q^{-1} + qx = qx - q^{-1}$. Dividing each by the normalization $q - q^{-1}x$ gives $\Ttilde_i(x)$ acting by $1$ on the $T_i = q$ eigenline and by $\rho(x) = (qx - q^{-1})/(q - q^{-1}x)$ on the $T_i = -q^{-1}$ eigenline. Multiplying numerator and denominator by $q$ yields $\rho(x) = (q^2 x - 1)/(q^2 - x)$.
\end{proof}

\begin{theorem}[M\"obius spectral bridge]
Define the M\"obius reparametrization
\[
\varphi(x) = \frac{q\big( (q^3 - 1) - (q-1) x \big)}{(q^3 - 1) x - q^2 (q - 1)}.
\]
Then
\[
\rho\big( \varphi(x) \big) = r(x) = \frac{q^2 - x}{q(x-1)}
\]
identically as rational functions in $q, x$. Equivalently, since both $\Rcheck_i$ and $\Ttilde_i$ have $P_q$-coefficient equal to $1$,
\[
\boxed{\ \Rcheck_i(x) = \Ttilde_i\big( \varphi(x) \big)\ }
\qquad \text{(exact equality of operators).}
\]
\end{theorem}

\begin{proof}
Both operators are $P_q + (\text{scalar}) \cdot P_{-1}$ in the same idempotent decomposition $I = P_q^{(i)} + P_{-1}^{(i)}$. The $P_q$-coefficients agree (both equal $1$). The $P_{-1}$-coefficients are $r(x)$ on the left and $\rho(\varphi(x))$ on the right; their equality is the displayed rational-function identity, verified symbolically over $\mathbb{Q}(q)$. Equality of the two scalar coefficients forces equality of the operators.
\end{proof}

\begin{corollary}[Special points]
$\varphi(q) = 1$, $\varphi(q^2) = q^{-2}$, and $\lim_{x \to 1} \varphi(x) = q^2$. Thus the dictionary of distinguished points is:
\begin{itemize}
  \item regular ($\Rcheck = I$): Clio $x = q$ $\leftrightarrow$ IK $x = 1$;
  \item fusion ($\Rcheck = P_q$): Clio $x = q^2$ $\leftrightarrow$ IK projector point $x = q^{-2}$;
  \item pole ($\Rcheck = P_{-1}$): Clio $x = 1$ $\leftrightarrow$ IK $x = q^2$.
\end{itemize}
\end{corollary}

\begin{remark}[Normalization invariance of the order]
If one uses the unnormalized IK $T_i(x)$ instead of $\Ttilde_i$, each factor is rescaled by the scalar
\[
f(x) = q - q^{-1} \varphi(x) = \frac{(q+1)(q^2+1)(x-1)}{q^2 x - q^2 + q x + x}.
\]
At the fusion point $f(q^2) = q - q^{-3}$, which is finite and nonzero. Hence the choice of normalized versus unnormalized IK generator multiplies $M(x)$ near $x = q^2$ by a finite nonzero scalar and does \emph{not} change $\ord_{x=q^2} Z_\lambda$. (All four identities above were verified symbolically over $\mathbb{Q}(q)$.)
\end{remark}

\section{Consequence: $\Omega$ lives in the Bethe construction}

Because $\Rcheck_i(x) = \Ttilde_i(\varphi(x))$ factor by factor, Clio's staircase $M(x)$ is, after the single global reparametrization $x \mapsto \varphi(x)$, an ordered staircase of genuine IK normalized Baxterised Hecke generators --- the same architecture as the IK Baxterised Jucys--Murphy element $y_n$. Therefore:

\begin{enumerate}
  \item[(i)] $Z_\lambda(x) = \tr_{V^\lambda} M(x)$ is a trace of a reparametrized IK transfer-type staircase: it lives in the Bethe-subalgebra \emph{trace} family.

  \item[(ii)] Clio's degenerate evaluation $x^* = q^2$ is, under $\varphi$, the IK projector/fusion point $x = q^{-2}$ where the normalized generator collapses to $P_q$ (the symmetric projector). So $\Omega$ is the fusion-point value of a traced staircase of IK Baxterised generators.

  \item[(iii)] Hence the order law $\ord_{x=q^2} Z_\lambda = \tau(\tau+1)/2$ is the order of vanishing of a Bethe transfer-matrix trace at a fusion point --- a $q$-KZ degeneration order in the Isaev--Kirillov--Tarasov (arXiv:1510.05374) sense, where these Baxterised elements furnish flat connections for the $q$-KZ equations.

  \item[(iv)] Clio's deficit law (deleting $u$ tail factors drops the $E^-$-sign-kill depth by exactly $u$) is then the concrete Hecke-basis computation of this degeneration's ramification: \emph{delete-$u$-drop-$u$} is \emph{merge-one-fewer-mode}.
\end{enumerate}

The pleasing thing is that nothing in this consequence is a new conjecture. Once the per-factor identity is exact, the architecture transports verbatim; what changes is only the interpretive frame in which $\Omega$ now sits.

\section{Honest scope and open questions}

\paragraph{What is proven here.}
The per-factor identity $\Rcheck_i(x) = \Ttilde_i(\varphi(x))$ and the special-point dictionary, exactly and symbolically. This places $M(x)$ and $Z_\lambda$ in the IK Baxterised-staircase / Bethe-trace family and identifies $x = q^2$ as the fusion point.

\paragraph{What is not yet established.}
\begin{enumerate}
  \item[(a)] The IK Bethe subalgebra is generated by the \emph{Markov} trace of $y_{n+1}$, whereas $Z_\lambda$ is a $V^\lambda$-character (an irreducible-module trace). Matching these requires expressing the $V^\lambda$-trace via Markov-trace coefficients --- a Geck--Rouquier / cocenter step.
  \item[(b)] The IK seed $y_1(x)$ and the inhomogeneities $\vec\xi$ have not been matched to Clio's homogeneous evaluation $\vec\xi = (1, \dots, 1)$, $x \to x$.
  \item[(c)] The quantitative count $\tau(\tau+1)/2$ is pinned to Clio's normalization; only the \emph{order} (not the leading coefficient) is shown normalization-invariant.
\end{enumerate}

Until (a)--(b) are settled, the correct statement is: $\Omega$ is structurally an IK Baxterised-staircase element evaluated at the fusion point, and a Bethe-subalgebra trace up to the trace-type identification.

\paragraph{Next step: the sharpened decisive test.}
Build the $(\tau+1) \times (\tau+1)$ confluent discrete Wronskian of the colliding antisymmetric-excess modes and verify that its ramification order is $\binom{\tau+1}{2}$ \emph{and} reproduces the deficit law's delete-$u$-drop-$u$ arithmetic. With the bridge in hand, this becomes an operator identity in $H_q(S_n)$ attackable with seminormal tools --- precisely the kind of statement that the spectral conjugacy was meant to make legible.

\end{document}
